\documentclass[10pt,a4paper]{amsart}
\usepackage[margin=19mm]{geometry}
\usepackage{amsmath,amssymb,mathtools}
\usepackage[scr=rsfs,cal=euler]{mathalfa}
\usepackage{xfrac}
\usepackage[unicode,hidelinks,bookmarks=false]{hyperref}
\newcommand{\ie}{i.e.\ }
\newcommand{\cf}{cf.\ }
\newcommand{\resp}{respectively\ }
\newcommand{\Subsection}{\subsection}
\providecommand{\nobreakdash}{\nobreak\hbox{-}}
\setlength{\emergencystretch}{2em}
\multlinegap0pt
\newtheorem{theorem}{Theorem}
\newtheorem{proposition}{Proposition}
\newtheorem{lemma}{Lemma}
\newtheorem{remark}{Remark}
\begin{document}
\title[An explicit construction of two completely independent spann]{An explicit construction of two completely independent spanning trees in the four-dimensional dual-cube}
\author{Jitendra Prajapati; Independent Researcher}
\date{}
\maketitle
\noindent\emph{Source and licence.} An explicit construction of two completely independent spanning trees in the four-dimensional dual-cube. Version arXiv:2608.00900v2 (CC BY 4.0). This material is distributed under the Creative Commons Attribution 4.0 International licence, http://creativecommons.org/licenses/by/4.0/, as recorded in the arXiv OAI licence metadata for this exact version and shown on the abstract page https://arxiv.org/abs/2608.00900v2. Full rights notice: licenses/arxiv-2608.00900-CC-BY-4.0.txt.\par\smallskip
\noindent\textit{Editorial scope.} Counted unit: the original \texttt{proposition} environment \texttt{prop:partition} at source lines 135--140 together with its complete proof at lines 142--146; the delivered document keeps source lines 53--186 so that every referenced label and every citation resolves inside it. Native editable source is the paper's own concatenated TeX; the AMS wrapper is editorial formatting only. No publisher class, upstream script or unrelated artwork is redistributed.
\begin{theorem}[\cite{Hasunuma2001}]\label{thm:hasunuma}
Spanning trees $T_1,\dots,T_k$ of a graph $G$ are completely independent if
and only if they are edge-disjoint and, for every vertex $x\in V(G)$, there
is at most one index $i$ with $d_{T_i}(x)>1$.
\end{theorem}

Hasunuma also proved that deciding the existence of two CISTs is
NP-complete~\cite{Hasunuma2002}. Araki~\cite{Araki2014} reformulated
existence in terms of vertex partitions; we use the following form, as
restated in~\cite{QinHaoWu2022}.

\begin{theorem}[\cite{Araki2014}]\label{thm:araki}
A connected graph $G$ has two completely independent spanning trees if and
only if $V(G)$ admits a partition $\{V_1,V_2\}$ such that $G[V_1]$ and
$G[V_2]$ are connected and the bipartite subgraph of $G$ induced by the
edge set $E(V_1,V_2)$ has no tree component. In that case the trees can be
chosen so that $V_i$ is exactly the set of internal vertices of $T_i$.
\end{theorem}

Constructions of two CISTs are known for hypercubes and several dense
hypercube variants~\cite{PaiChang2016}; see the survey~\cite{ChengWangFan2023}.
The dual-cube $F_n$, introduced by Li and Peng~\cite{LiPeng2000}, is sparse
by design, which places it near the edge-count threshold for two
edge-disjoint spanning trees. Lalou, Mbarek, Skender and
Togni~\cite{Lalou2026} recently proved that $F_n$ admits two CISTs for
every $n\ge 5$, observed that the edge count already obstructs existence
for $n\le 3$, and reported that the case $n=4$ resisted more than $700$
hours of computation and was left unresolved.

In this note we settle the remaining case.

\begin{theorem}\label{thm:main}
$F_4$ admits two completely independent spanning trees. Consequently, the
dual-cube $F_n$ admits two completely independent spanning trees if and
only if $n\ge 4$.
\end{theorem}

The construction (Section~\ref{sec:construction}) is compact: the
internal-vertex partition is given by a ten-term cubic polynomial over
$\mathbb{F}_2$ in the seven vertex bits, and its correctness reduces via
Theorem~\ref{thm:araki} to finite connectivity checks. A certificate
pinning the two trees explicitly, together with a verifier that uses only
the Python standard library and checks all $\binom{128}{2}=8128$ pairs of
tree paths directly against the definition, is publicly available
(Section~\ref{sec:verification}). Section~\ref{sec:searches} reports exact
infeasibility results for simpler rules of the same shape.

\section{Preliminaries}\label{sec:prelim}

Following~\cite{Lalou2026}, the dual-cube $F_n$ is the $n$-regular
$n$-connected graph on the vertex set $\{0,1\}^{2n-1}$ in which two
vertices $u=(u_{2n-1},\dots,u_1)$ and $v=(v_{2n-1},\dots,v_1)$ are adjacent
if and only if they differ in exactly one bit position $i$ and:
\begin{enumerate}
\item if $1\le i\le n-1$ then $u_{2n-1}=v_{2n-1}=0$;
\item if $n\le i\le 2n-2$ then $u_{2n-1}=v_{2n-1}=1$;
\item $i=2n-1$ is always allowed.
\end{enumerate}
Thus $F_4$ has $128$ vertices, $256$ edges, and degree $4$; it consists of
sixteen $Q_3$ clusters, eight in each class, joined by a perfect matching
of $128$ cross-edges.

For consistency with the machine-readable certificate we index bits from
zero: a vertex $v\in\{0,1\}^7$ of $F_4$ has bits $x_0,\dots,x_6$ with
$x_i=u_{i+1}$, so $x_6$ is the class bit, $x_0,x_1,x_2$ are the class-$0$
cluster coordinates and $x_3,x_4,x_5$ are the class-$1$ cluster
coordinates.

Two edge-disjoint spanning trees of $F_4$ use $2\cdot 127=254$ of the $256$
edges. By Theorem~\ref{thm:hasunuma}, in any CIST pair no vertex is
internal in both trees; since $F_4$ is $4$-regular, every vertex is then
either internal in exactly one tree or a leaf in both.

\section{The construction}\label{sec:construction}

Define $z\colon V(F_4)\to\mathbb{F}_2$ by
\begin{equation}\label{eq:anf}
z(v) \;=\; x_0x_1 \oplus x_2 \oplus x_1x_2 \oplus x_1x_3 \oplus x_3x_4
\oplus x_1x_5 \oplus x_4x_5 \oplus x_0x_6 \oplus x_3x_6 \oplus x_2x_4x_6,
\end{equation}
and set $V_1=\{v : z(v)=1\}$ and $V_2=\{v : z(v)=0\}$.


\begin{proposition}\label{prop:partition}
$|V_1|=|V_2|=64$; the induced subgraphs $F_4[V_1]$ and $F_4[V_2]$ are
connected with exactly $64$ edges each (hence unicyclic); and the bipartite
subgraph induced by the $128$ edges of $E(V_1,V_2)$ has exactly two
components, each with $64$ vertices and $64$ edges (hence unicyclic).
\end{proposition}

\begin{proof}
This is a finite computation on a $128$-vertex graph; it is performed
independently of any solver by the verification program of
Section~\ref{sec:verification}.
\end{proof}


Since a unicyclic component is not a tree, Proposition~\ref{prop:partition}
and Theorem~\ref{thm:araki} prove Theorem~\ref{thm:main}: $\{V_1,V_2\}$ is
a CIST partition of $F_4$, and the classification follows
from~\cite{Lalou2026}. Every vertex is internal in exactly one tree, so the
pair realizes the balanced $64/64$ pattern.

The certificate pins the two trees explicitly, so correctness does not
depend on quoting Theorem~\ref{thm:araki}. The unique cycle of $F_4[V_2]$
contains the edge $\{0,64\}$ and the unique cycle of $F_4[V_1]$ contains
the edge $\{21,85\}$. Orient each unicyclic component of the cut graph
functionally, so that every vertex acquires exactly one outgoing cut edge
$\{u,\mathrm{out}(u)\}$. Then
\[
T_2^{\,\prime} = \bigl(E(F_4[V_2])\setminus\{\{0,64\}\}\bigr)\;\cup\;
\bigl\{\{u,\mathrm{out}(u)\} : u\in V_1\bigr\},
\]
and symmetrically $T_1^{\,\prime}$ with $\{21,85\}$ removed, are two
edge-disjoint spanning trees, with internal-vertex sets exactly $V_1$ and
$V_2$; the only edges of $F_4$ unused by $T_1^{\,\prime}\cup
T_2^{\,\prime}$ are $\{0,64\}$ and $\{21,85\}$.

\section{Machine verification}\label{sec:verification}

The construction is distributed as a short JSON certificate together with a
verifier that uses only the Python standard library --- no solver is
involved. The verifier (i)~rebuilds $F_4$ from the bit definition of
Section~\ref{sec:prelim}, (ii)~evaluates the polynomial~\eqref{eq:anf} and
checks every statement of Proposition~\ref{prop:partition},
(iii)~reconstructs $T_1^{\,\prime}$ and $T_2^{\,\prime}$ and checks that
they are edge-disjoint spanning trees with internal-vertex sets $V_1$ and
$V_2$, and (iv)~for each of the $8128$ unordered vertex pairs checks
directly that the two tree paths are openly disjoint. Step~(iv) is the
definition of complete independence, so the verified claim does not rest on
Theorem~\ref{thm:araki} or on any search software. A second, independently
fixed edge-list witness passes a separate verifier. Both programs refuse to
run under Python's assertion-stripping \texttt{-O} mode.

\section{Exact searches for simpler rules}\label{sec:searches}


\begin{thebibliography}{99}
\bibitem[Araki2014]{Araki2014} T.~Araki, Dirac's condition for completely independent spanning trees, \emph{J. Graph Theory} 77 (2014) 171--179.
\bibitem[ChengWangFan2023]{ChengWangFan2023} B.~Cheng, D.~Wang, and J.~Fan, Independent spanning trees in networks: a survey, \emph{ACM Comput. Surv.} 55(14s) (2023) Article 335.
\bibitem[Hasunuma2001]{Hasunuma2001} T.~Hasunuma, Completely independent spanning trees in the underlying graph of a line digraph, \emph{Discrete Math.} 234 (2001) 149--157.
\bibitem[Hasunuma2002]{Hasunuma2002} T.~Hasunuma, Completely independent spanning trees in maximal planar graphs, in: \emph{Proc. 28th Int. Workshop on Graph-Theoretic Concepts in Computer Science (WG 2002)}, LNCS 2573, Springer, 2002, pp. 235--245.
\bibitem[Lalou2026]{Lalou2026} M.~Lalou, N.~Mbarek, A.~Skender, and O.~Togni, Constructing two completely independent spanning trees in the dual-cube, arXiv:2607.25917 (2026).
\bibitem[LiPeng2000]{LiPeng2000} Y.~Li and S.~Peng, Dual-cubes: a new interconnection network for high-performance computer clusters, in: \emph{Proc. Int. Computer Symposium}, 2000.
\bibitem[PaiChang2016]{PaiChang2016} K.-J.~Pai and J.-M.~Chang, Constructing two completely independent spanning trees in hypercube-variant networks, \emph{Theor. Comput. Sci.} 652 (2016) 28--37.
\bibitem[QinHaoWu2022]{QinHaoWu2022} X.-W.~Qin, R.-X.~Hao, and J.~Wu, Construction of dual-CISTs on an infinite class of networks, \emph{IEEE Trans. Parallel Distrib. Syst.} 33(8) (2022) 1902--1910.
\end{thebibliography}
\end{document}
